% Created 2026-09-16 Wed 13:40
% Intended LaTeX compiler: pdflatex
\documentclass[presentation]{beamer}
\usepackage[utf8]{inputenc}
\usepackage[T1]{fontenc}
\usepackage{graphicx}
\usepackage{longtable}
\usepackage{wrapfig}
\usepackage{rotating}
\usepackage[normalem]{ulem}
\usepackage{amsmath}
\usepackage{amssymb}
\usepackage{capt-of}
\usepackage{hyperref}
\usetheme{default}
\author{Rick Wysocki}
\date{}
\title{Week 4 / Thursday}
\hypersetup{
 pdfauthor={Rick Wysocki},
 pdftitle={Week 4 / Thursday},
 pdfkeywords={},
 pdfsubject={},
 pdfcreator={Emacs 30.2 (Org mode 9.7.11)}, 
 pdflang={English}}
\begin{document}

\maketitle
\begin{frame}[label={sec:orgd3a910a}]{Markup Languages}
A markup language, which is \alert{distinct from "writing code,"} can be defined as:

\begin{quote}
a text-encoding system which specifies the structure and formatting of a document and potentially the relationships among its parts. Markup can control the display of a document or enrich its content to facilitate automated processing.

-- \href{https://en.wikipedia.org/wiki/Markup\_language}{Wikipedia}
\end{quote}
\end{frame}
\begin{frame}[label={sec:orgacfffbc}]{WYSIWYG vs. Plain Text Editors}
WYSIWYG ("What You See is What You Get") editors display formatted text while you type, whereas plain text editors omit formatting, often so that the "plain" text can be put into alternative displays.
\end{frame}
\begin{frame}[label={sec:org2c0eab4}]{Why Plain Text?}
Plain text writing separates the domains of content and presentation (also sometimes referred to as views or displays) in digital environments. For what reasons might this be useful?
\end{frame}
\begin{frame}[label={sec:org6dec755}]{Examples of Plain Text Markup Languages}
\begin{itemize}
\item HTML
\item XML
\item Tex (LaTex)
\item Markdown
\end{itemize}
\end{frame}
\begin{frame}[label={sec:org22b6ade}]{Key Affordances}
\begin{itemize}
\item Separating content from presentation
\item Modularity
\item Reuse
\item Stability and preservation
\item Accessibility
\item Intentionality
\end{itemize}
\end{frame}
\end{document}
